% Cierre: material de partida y fuentes originales
\phantomsection\section*{Material de partida}\label{cap:22}
\addcontentsline{toc}{section}{Material de partida}

\begin{itemize}
  \item \texttt{Ejercicios.pdf}: definiciones, Teoremas 1 a 6 y consigna de
  verificación geométrica y algebraica.
  \item Material sobre la primera sesión de rotaciones: rotaciones 2D y 3D, números
  complejos, orden de composición y exponencial de matrices.
  \item \texttt{02.afinidad.pdf}: repaso de rotaciones, composición de
  reflexiones, puntos fijos, forma afín y coordenadas homogéneas.
  \item Material práctico de rotación de \(45^\circ\),
  escala \(0.5\), composición matricial, Excel y punto fijo.
  \item Material sobre el plano afín \(w=1\),
  rayos homogéneos, puntos ideales y normalización proyectiva.
  \item Nota sobre transformaciones proyectivas: ``mantienen rectas las rectas
  \(\Rightarrow\) quedan definidas por un cuadrilátero''.
  \item Material sobre producto interno e
  incidencia, producto vectorial, determinante, producto exterior de Grassmann, núcleo,
  autovalores y autovectores por producto cruz de filas, ejemplos de reflexión y rotación,
  y consigna del ejercicio del pentágono regular.
  \item Programa oficial y resumen conceptual existente: coordenadas curvilíneas,
  variedades, métrica, elipsoide, proyecciones, marcos y grillas.
  \item Resúmenes conceptuales de las diez unidades del esqueleto de la materia, con sus
  indicaciones bibliográficas.
  \item Bibliografía comentada (apéndice~\ref{ap:bibliografia}) y material práctico
  (planillas, notebooks de wxMaxima, ejemplo en Python).
  \item Apunte sobre el elipsoide, ``Geometría del Elipsoide y cálculos Geodésicos'' (guía
  de 20 páginas): fuente de las
  secciones~\ref{cap:curvatura} a \ref{cap:esfera}.
  \item Planilla de la proyectiva del pentágono y cuaderno \file{LaSVD.wxmx}:
  apéndice~\ref{ap:notebooks}.
  \item Fichas de unidades y unidades interactivas del curso (\file{clases/}, \file{unidades/} del
  repositorio).
\end{itemize}

\vfill
\begin{center}
\color{muted}\small
Documento preparado para estudio. Conviene rehacer a mano cada figura y cada cuenta:
la comprensión está en la construcción, no solamente en el resultado.
\end{center}

\section*{Fuentes originales}
\addcontentsline{toc}{section}{Fuentes originales}
Las páginas siguientes reproducen íntegramente los documentos de partida recibidos, en el
orden en que fueron enviados:
\begin{enumerate}
  \item guía de definiciones, teoremas y ejercicios;
  \item resumen sobre transformaciones afines.
\end{enumerate}
